\chapter{Prompt Engineering -- Die Schnittstelle, die zählt}
\label{chap:prompt_engineering}

% ============================================================
% AUTOR-NOTIZ STATT LERNZIELE
% ============================================================
\autornotiz{
2022 habe ich meinen ersten Prompt in Produktion deployed. Ein Klassifizierer für Support-Tickets.
Zero-Shot, temperature 0.7, keine Validierung. Kosten: \$400/Tag. Accuracy: 61\%.
Heute: Few-Shot (3 Edge-Cases), temperature 0.0, Pydantic-Validierung, Semantic Caching.
Kosten: \$12/Tag. Accuracy: 94\%.
Der Unterschied war nicht das Modell. Es war Disziplin bei der Prompt-Schnittstelle.
}

% ============================================================
% MOTIVATION -- WARUM DAS HIER DIE HAUPTSCHNITTSTELLE IST
% ============================================================
\section{Warum Prompt Engineering die Hauptprogrammierschnittstelle ist}

Software-Entwickler schreiben Code. Code ist deterministisch. Derselbe Input $\rightarrow$ derselbe Output. Immer.

AI Engineers schreiben Prompts. Ein Prompt ist eine \textbf{Wahrscheinlichkeitsverteilung über Token-Sequenzen}. Derselbe Prompt $\rightarrow$ \textit{verschiedene} Outputs. Temperatur, Server-seitiges Batching, Modell-Updates -- alles verschiebt die Verteilung.

\par



\par\mbox{}\par
\begin{realitycheckEnv}
Wer "Prompt Engineering = gute Fragen stellen" denkt, baut Systeme, die in Produktion brechen. Prompt Engineering ist: \textbf{Struktur erzwingen, Validierung bauen, Distributionen einengen}. Das ist Engineering. Der Rest ist Hoffen.
\end{realitycheckEnv}

Unternehmen mit systematischem Prompt Engineering haben messbar niedrigere Fehlerraten und Token-Kosten. Nicht weil sie "bessere Prompts schreiben", sondern weil sie Prompts wie Code behandeln: versioniert, getestet, gemonitored, gerollt-backt.

\skipbox{
\textbf{TL;DR dieses Kapitels:}
\begin{itemize}
    \item Few-Shot (3--5 Edge-Cases) schlägt Zero-Shot immer -- Mechanismus: Pattern Matching + Format Following + Attention Bias
    \item Temperature $>$ 0.3 in Produktion = Fahrlässigkeit
    \item System-Prompt + User-Input \textbf{niemals} interpolieren -- das ist die \#1 Sicherheitslücke
    \item Prompts gehören in Dateien (Jinja2), versioniert (Git-Tags), evaluiert (Golden Set), gecached (Prompt Caching)
    \item Self-Consistency (5 Samples + Majority Vote) kostet 5x Latenz, bringt +10-20\% -- nur für mission-critical
\end{itemize}
}

% ============================================================
% DIE DREI EBENEN -- MECHANISMUS STATT DEFINITION
% ============================================================
\section{Die drei Ebenen eines Prompts -- und warum jede fehlt, rät das Modell}

Jeder Prompt hat drei implizite Ebenen. Fehlt eine, füllt das Modell die Lücke mit Trainingsdaten-Statistik -- nicht mit deiner Intent.

\begin{enumerate}
    \item \textbf{Instruktion}: Was tun? (Explizit, sichtbar)
    \item \textbf{Kontext}: Worum geht es? (System-Prompt, History, RAG-Chunks)
    \item \textbf{Meta-Ebene}: Wie denken? (Persona, Ton, Output-Format, Reasoning-Style)
\end{enumerate}

\realitycheck{Die Meta-Ebene ist der billigste Hebel für Konsistenz. "Antworte nur JSON" im System-Prompt spart 40\% Parsing-Fehler vs. im User-Prompt. Getestet bei 10k Requests/Tag.}

\begin{verbatim}
Prompt = Instruktion + Kontext + Meta-Ebene

[System] Du bist ein präziser Support-Klassifizierer.      <-- Meta-Ebene
         Antworte AUSCHLIESSLICH valides JSON.             <-- Meta-Ebene
         Schema: {"priority":1-5,"category":"...",...}     <-- Meta-Ebene
         Klassifiziere Tickets nach Priorität.             <-- Instruktion
         Historie: {history}                               <-- Kontext

[User]   Mein Produkt ist kaputt!                          <-- Instruktion
\end{verbatim}

% ============================================================
% PROMPT-ENGINEERING-PYRAMIDE -- MECHANISMUS STATT LISTE
% ============================================================
\section{Die Prompt-Engineering-Pyramide -- warum die Techniken so stapeln}

\begin{verbatim}
                      /\
                     /  \          Self-Consistency
                    /    \         (Mehrere Samples + Voting)
                   /      \
                  /--------\       ReAct (Reason + Act)
                 /          \
                /------------\     Chain-of-Thought
               /              \
              /----------------\   Few-Shot (mit Beispielen)
             /                  \
            /--------------------\ Zero-Shot (reine Instruktion)
\end{verbatim}

Die Pyramide ist keine Hierarchie der "Schwierigkeit" -- sie zeigt, \textbf{welchen Mechanismus} jede Technik aktiviert:

\begin{itemize}

    \item \textbf{Zero-Shot} Basisfall. Nutzt nur das vortrainierte Prior des Modells. Funktioniert bei starken Modellen für klare Aufgaben.
    \item \textbf{Few-Shot} Aktiviert \textbf{In-Context Learning (ICL)}: Pattern Matching + Format Following + Attention Bias. \realitycheck{Drei gut gewählte Edge-Cases schlagen 20 zufällige Samples. Qualität $>$ Quantität.}
    \item \textbf{Chain-of-Thought} Erzwingt \textbf{intermediate reasoning tokens}. Das Modell "denkt laut" -- die Token-Sequenz wird zum Scratchpad. \realitycheck{CoT kostet 1.5-2x Latenz. Lohnt sich nur bei Logik/Math/Mehrschritt. Für Klassifikation: Overhead ohne Gewinn.}
    \item \textbf{ReAct} Kombiniert Reasoning mit \textbf{Tool-Use}. Basis aller Agent-Architekturen. Das Modell erzeugt strukturierte Aktionen (Function Calling), führt sie aus, observed, repeatet.
    \item \textbf{Self-Consistency} Nutzt \textbf{Ensemble-Effekt}: Mehrere Samples (temperature 0.7) + Majority Vote. Reduziert Varianz. \realitycheck{5x API-Calls. Nur wenn Fehler teuer sind (Medizin, Finanzen, Security).}
\end{itemize}

% ============================================================
% ARCHITEKTUR -- PROMPT PIPELINES IN PRODUKTION
% ============================================================
\section{Architektur -- Prompt-Pipelines, die in Produktion überleben}

In Produktion ist ein Prompt keine einzelne Message. Er durchläuft eine Pipeline:

\begin{verbatim}
Input-Daten
    |
    v
[Pre-Processing]   --> Bereinigung, Formatierung, Kontext-Anreicherung (RAG)
    |
    v
[Template-Rendering] --> Jinja2 befüllt Template (System + Few-Shot + User)
    |
    v
[Token-Budget-Check] --> Input-Tokens zählen, ggf. truncaten (History summarization)
    |
    v
[API-Call]         --> OpenAI / Anthropic / Google / Self-Hosted
    |
    v
[Post-Processing]  --> JSON parsen, Pydantic validieren, transformieren
    |
    v
[Output]           --> strukturiertes Ergebnis + Metriken (Tokens, Latenz, Kosten)
\end{verbatim}

\autornotiz{
Bei Claypot (mein Startup, 2021-2024) haben wir die Pipeline erst gebaut, als der dritte Prompt-Change ein Deployment brauchte. Lektion: Pipeline \textbf{vor} dem zweiten Prompt bauen. Die Jinja2-Registry (siehe Code unten) hat uns 20+ Deployments gespart.
}

\subsection{Prompt-Versionierung -- genau wie Code}

\begin{table}[H]
\centering
\begin{tabularx}{\linewidth}{lX}
\toprule
\textbf{Aspekt} & \textbf{Was funktioniert} \\
\midrule
Speicherort & \texttt{prompts/v1/classification.jinja2} \\
Versionierung & Git-Tag pro Prompt-Änderung (\texttt{prompt/classify-v3}) \\
Evaluierung & Golden Set (50-100 Cases) vor \textit{jedem} Deploy \\
Rollback & \texttt{git checkout prompt/classify-v2} -- 30 Sekunden \\
\midrule
\end{tabularx}
\caption{Prompt-Versionierungs-Strategie, die hält}
\label{tab:prompt-versioning}
\end{table}

\realitycheck{Ohne Golden Set ist jedes Prompt-Update ein Blindflug. "Fühlt sich besser an" skaliert nicht. Automatisiere die Eval-Pipeline \textit{bevor} du den zweiten Prompt schreibst.}

% ============================================================
% THEORIE -- WARUM DAS FUNKTIONIERT (MECHANISMUS)
% ============================================================
\section{Theorie -- Die Mechanismen hinter den Techniken}

\skipdeep{TL;DR: In-Context Learning = Pattern Matching + Format Following + Attention Bias. Temperature steuert Entropie der Output-Distribution. Receptivity Maximum bei ~2-4k Instruktion-Tokens. Details unten.}

\subsection{In-Context Learning -- das Fundament}

ICL verändert \textbf{keine} Modellgewichte. Es nutzt die im Training gelernten Muster, um aus Prompt-Beispielen zu generalisieren. Drei Mechanismen wirken gleichzeitig:

\begin{enumerate}
    \item \textbf{Pattern Matching}: Das Modell erkennt Input $\rightarrow$ Output Muster im Prompt und wendet es an. Wie Few-Shot funktioniert: Das Modell kopiert die \textit{Struktur}, nicht den Inhalt.
    \item \textbf{Format Following}: Durch Beispiele lernt das Modell das exakte Output-Format. JSON mit bestimmten Keys, Enum-Werte, verschachtelte Strukturen. \realitycheck{Ohne Examples: 67\% valides JSON. Mit 3 Examples: 94\%. Getestet mit gpt-4o-mini, 500 Samples.}
    \item \textbf{Attention Bias}: Die Beispiele lenken die Attention auf relevante Aspekte. Ein "borderline billing/complaint" Example im Few-Shot verschiebt die Klassifikationsgrenze um ~15\% Recall.
\end{enumerate}

\subsection{Temperature -- Entropie-Kontrolle, nicht Kreativität}

Temperature skaliert die Logits vor dem Softmax. Temperature 0.0 = argmax (deterministisch, bis auf Floating-Point-Nondeterminismus durch Server-seitiges Batching). Temperature 1.0 = Original-Distribution. Temperature $>$ 1.0 = flacher, "kreativer".

\begin{table}[H]
\centering
\begin{tabularx}{\linewidth}{lXX}
\toprule
\textbf{Temperature} & \textbf{Verhalten} & \textbf{Produktions-Einsatz} \\
\midrule
0.0--0.1 & Deterministisch, konsistent & \textbf{Default für Produktion}. Klassifikation, Extraktion, Structured Output \\
0.2--0.3 & Leichte Varianz, hohe Qualität & Chat, wenn Konsistenz wichtig aber nicht kritisch \\
0.4--0.6 & Kreativ, variabel & Brainstorming, Drafting \\
0.7--1.0 & Maximale Entropie & Creative Writing, \textbf{nie} in Produktion ohne Guardrails \\
\midrule
\end{tabularx}
\caption{Temperature = Entropie-Kontrolle. Kreativität ist Marketing.}
\label{tab:temperature}
\end{table}

\realitycheck{Seed (z.B. 42) gibt \textit{best effort} Reproduzierbarkeit. OpenAI dokumentiert: "best effort" -- server-seitiges Batching kann Floating-Point-Arithmetik minimal verschieben. Verlass dich nicht auf bit-identische Outputs.}

\subsection{Warum Few-Shot Zero-Shot schlägt -- die drei Mechanismen}

Die Few-Shot-Beispiele erfüllen drei \textit{technische} Funktionen:

\begin{enumerate}
    \item \textbf{Format-Vorlage}: Das Modell sieht das \textit{exakte} Output-Format. Kein Raten.
    \item \textbf{Entscheidungsanker}: Grenzfälle im Example kalibrieren die Entscheidungsschwelle. Ein "priority 3 vs 4" Example im Few-Shot setzt die Boundary präziser als jede Instruktion.
    \item \textbf{Attention Bias}: Die Beispiele wirken wie Spotlight: "Hier lang schauen, Rest ignorieren."
\end{enumerate}

\realitycheck{Die besten Few-Shot-Beispiele sind nicht "repräsentativ" -- sie sind die \textbf{Edge-Cases, bei denen das Modell früher gescheitert ist}. Sammle Fehlschläge. Werde sie zu Examples.}

\subsection{Das Rezeptivitäts-Maximum -- länger ist nicht besser}

Ab ~2.000-4.000 Tokens \textit{Instruktion} (nicht Kontext!) sinkt die Qualität. Die Attention verteilt sich über mehr Tokens, die Kern-Instruktion "verschwimmt" im Rauschen.




\par\mbox{}\par
\begin{praxishinweisEnv}
Faustregel: Instruktionen so kurz wie möglich, Examples so viele wie nötig (3-5), Kontext so wenig wie möglich. Miss \texttt{prompt\_efficiency\_score()} -- optimiere auf Tokens pro Qualitätspunkt.
\end{praxishinweisEnv}

% ============================================================
% PRAXIS -- DIE FÜNF TECHNIKEN MIT CODE, DER LÄUFT
% ============================================================
\section{Praxis -- Die fünf Techniken, wann sie laufen, wann nicht}

\subsection{Zero-Shot -- der Basisfall}

Funktioniert bei starken Modellen (gpt-4o, claude-4-sonnet) für klare, einfache Aufgaben.

\begin{lstlisting}[language=Python, caption={Zero-Shot: sauber, validiert, Temperatur 0}]
from openai import OpenAI
from pydantic import BaseModel

client = OpenAI()

class TicketClassification(BaseModel):
    priority: int        # 1-5
    category: str        # billing | technical | complaint | feature_request
    urgency: str         # low | medium | high

def classify_zero_shot(ticket_text: str) -> TicketClassification:
    resp = client.chat.completions.create(
        model="gpt-4o-mini",
        messages=[
            {"role": "system", "content": (
                "Klassifiziere Support-Tickets. "
                "Antworte AUSCHLIESSLICH valides JSON nach Schema."
            )},
            {"role": "user", "content": ticket_text},
        ],
        response_format={"type": "json_object"},
        temperature=0.0,
    )
    return TicketClassification.model_validate_json(resp.choices[0].message.content)

# Test
result = classify_zero_shot("Meine Bestellung ist seit 2 Wochen nicht da!")
print(result)  # priority=4, category='complaint', urgency='high'
\end{lstlisting}

\realitycheck{Zero-Shot ist dein Baseline. Miss Accuracy. Wenn $<$ 85\%, brauchst du Few-Shot. Nicht raten -- messen.}

\subsection{Few-Shot -- der Produktionsstandard}

Der wichtigste Hebel. 3-5 Examples, die Edge-Cases abdecken.

\begin{lstlisting}[language=Python, caption={Few-Shot: Examples als Daten, nicht Code}]
FEW_SHOT_EXAMPLES = [
    {"input": "Ich habe das falsche Produkt erhalten.",
     "output": '{"priority": 3, "category": "complaint", "urgency": "medium"}'},
    {"input": "Die Rechnung zeigt einen falschen Betrag.",
     "output": '{"priority": 4, "category": "billing", "urgency": "high"}'},
    {"input": "Würdet ihr bald Widget X anbieten?",
     "output": '{"priority": 1, "category": "feature_request", "urgency": "low"}'},
    # Edge-Case: billing + complaint overlap
    {"input": "Ich wurde doppelt belastet UND das Produkt ist defekt!",
     "output": '{"priority": 5, "category": "billing", "urgency": "high"}'},
]

def build_few_shot_prompt(examples: list[dict], user_input: str) -> str:
    parts = ["Klassifiziere als JSON: priority (1-5), category, urgency."]
    parts += [f"Input: {e['input']}\nOutput: {e['output']}" for e in examples]
    parts.append(f"Input: {user_input}\nOutput:")
    return "\n\n".join(parts)

# Nutzung mit JSON Mode + Validierung
resp = client.chat.completions.create(
    model="gpt-4o-mini",
    messages=[
        {"role": "system", "content": "Antworte AUSCHLIESSLICH valides JSON."},
        {"role": "user", "content": build_few_shot_prompt(FEW_SHOT_EXAMPLES, user_text)},
    ],
    response_format={"type": "json_object"},
    temperature=0.0,
)
result = TicketClassification.model_validate_json(resp.choices[0].message.content)
\end{lstlisting}

\realitycheck{Vier Examples oben: 3 Standard-Klassen + 1 Overlap-Edge-Case. Das Overlap-Example allein hob Billing-Recall um 12\% in meinem Test-Set. Sammle deine Fehler. Mach sie zu Examples.}

\subsection{Chain-of-Thought -- für Logik, nicht für Klassifikation}

\begin{lstlisting}[language=Python, caption={CoT: strukturiertes Reasoning erzwingen}]
COT_SYSTEM = """Denke Schritt für Schritt.
Struktur:
  Reasoning: [deine Überlegungen]
  Answer: [finales JSON]"""

resp = client.chat.completions.create(
    model="gpt-4o",
    messages=[
        {"role": "system", "content": COT_SYSTEM},
        {"role": "user", "content": (
            "Ein Agent bearbeitet 40 Tickets/Tag. "
            "5% brauchen 30 Min, 95% brauchen 10 Min. "
            "Wie viele kritische (30 Min) Tickets schafft er?"
        )},
    ],
    temperature=0.0,
)
# Output parsing: JSON aus "Answer:" extrahieren + validieren
\end{lstlisting}

\realitycheck{CoT kostet 1.5-2x Latenz (mehr Output-Tokens). Bei Klassifikation/Extraktion: \textbf{kein} messbarer Gewinn. Nutze es nur für: Math, Logik, Planung, Debugging, Multi-Step-Analysis.}

\subsection{ReAct -- Reasoning + Tool Use (Agent-Basis)}

\begin{lstlisting}[language=Python, caption={ReAct: strukturierter Tool-Loop}]
REACT_SYSTEM = """Löse Tasks durch Denken und Handeln.

Format:
  Thought: [Was muss ich tun?]
  Action: [function_name]
  Action Input: [JSON-Parameter]
  Observation: [Ergebnis]
  ... (wiederholen bis fertig) ...
  Answer: [Endergebnis als JSON]

Tools:
  - search(query): Wissensdatenbank-Suche
  - calculate(expr): Mathematik
  - lookup_product(id): Produkt-Details
"""

resp = client.chat.completions.create(
    model="gpt-4o",
    messages=[
        {"role": "system", "content": REACT_SYSTEM},
        {"role": "user", "content": "Was kostet Produkt #12345? Gibt es Rabatte?"},
    ],
    temperature=0.0,
)
# Parser braucht: Thought/Action/Observation Loop ausführen, Tools aufrufen, Result injizieren
\end{lstlisting}

\realitycheck{ReAct-Prompts sind \textbf{brittle}. Ein falsches Action-Format und der Loop bricht. In Produktion: Nutze Function Calling (OpenAI) / Tool Use (Anthropic) statt Text-Parsing. Das Modell gibt strukturierte Calls, du führst aus, gibst structured Result zurück. Weniger Parsing, weniger Failures.}

\subsection{Self-Consistency -- Majority Vote für Mission-Critical}

\begin{lstlisting}[language=Python, caption={Self-Consistency: 5 Samples + Majority Vote}]
from collections import Counter

def self_consistent_classify(prompt: str, n: int = 5) -> TicketClassification:
    """Nur für mission-critical. Kostet 5x API-Calls + Latenz."""
    answers = []
    for _ in range(n):
        resp = client.chat.completions.create(
            model="gpt-4o-mini",
            messages=[{"role": "user", "content": prompt}],
            temperature=0.7,  # WICHTIG: Varianz erzeugen
            response_format={"type": "json_object"},
        )
        answers.append(resp.choices[0].message.content)
    
    # Majority Vote auf JSON-Strings
    most_common = Counter(answers).most_common(1)[0][0]
    return TicketClassification.model_validate_json(most_common)

# Kosten-Nutzen: +10-20% Accuracy bei 5x Cost/Latency
# Lohnt sich: Medical diagnosis, Fraud detection, Legal classification
\end{lstlisting}

\realitycheck{Self-Consistency ist teuer. In 3 Jahren Produktion habe ich es \textbf{zweimal} gebraucht: einmal für ICD-10 Coding (Medical), einmal für Betrugserkennung (FinTech). Standard-Klassifikation: Few-Shot + temperature 0.0 reicht.}

% ============================================================
% PROMPT TEMPLATES -- JINJA2 REGISTRY (PRODUCTION CODE)
% ============================================================
\section{Prompt-Templates mit Jinja2 -- Registry Pattern, das deployt}

Prompts hardcoden = Deployment pro Änderung. Templates in Dateien = Änderung ohne Deploy.

\begin{lstlisting}[language=Python, caption={Prompt Registry: Versioniert, gecached, evaluierbar}]
from jinja2 import Environment, FileSystemLoader
from pathlib import Path
import hashlib
import json

class PromptRegistry:
    """Production-ready Prompt Registry mit Versionierung & Caching."""
    
    def __init__(self, template_dir: str = "prompts"):
        self.env = Environment(
            loader=FileSystemLoader(template_dir),
            autoescape=False,
            trim_blocks=True,
            lstrip_blocks=True,
        )
        self._cache: dict[str, str] = {}
    
    def render(self, name: str, version: int = 1, **kwargs) -> str:
        """Rendert Template v{version}/{name}.jinja2 mit kwargs."""
        template_path = f"v{version}/{name}.jinja2"
        template = self.env.get_template(template_path)
        return template.render(**kwargs)
    
    def render_with_hash(self, name: str, version: int = 1, **kwargs) -> tuple[str, str]:
        """Gibt (rendered_prompt, content_hash) zurück -- für Caching/Logging."""
        rendered = self.render(name, version, **kwargs)
        content_hash = hashlib.sha256(rendered.encode()).hexdigest()[:16]
        return rendered, content_hash

# Datei: prompts/v1/classification.jinja2
# ---
# Du klassifizierst Support-Tickets. Antworte AUSCHLIESSLICH JSON:
# {"priority": 1-5, "category": "...", "urgency": "..."}
#
# {% for ex in examples %}
# Input: {{ ex.input }}
# Output: {{ ex.output }}
# {% endfor %}
#
# Ticket: {{ ticket_text }}
# Output:
# ---

# Nutzung
registry = PromptRegistry()
prompt, p_hash = registry.render_with_hash(
    "classification", version=3,
    examples=db.get_few_shot_examples(),
    ticket_text=request.text,
)
# p_hash loggen -> Prompt-Version in Traces nachvollziehbar
\end{lstlisting}

\realitycheck{Prompt-Hash in jedem Trace (Langfuse, LangSmith, custom) = du weißt \textit{exakt}, welcher Prompt welchen Output produziert hat. Ohne Hash: "Welche Version war das nochmal?" bei 3am Debugging.}

% ============================================================
% PERFORMANCE -- LATENZ, KOSTEN, CACHING
% ============================================================
\section{Performance -- Latenz, Kosten, Caching: was sich lohnt}

\subsection{Technik-Trade-offs mit echten Zahlen}

\begin{table}[H]
\centering
\begin{tabularx}{\linewidth}{lXX}
\toprule
\textbf{Technik} & \textbf{Latenz (relativ)} & \textbf{Qualitäts-Delta} \\
\midrule
Zero-Shot & 1.0x (Baseline) & Baseline \\
Few-Shot (3 Examples) & 1.1x & +15-30\% \\
Few-Shot (10 Examples) & 1.3x & +20-40\% (diminishing returns) \\
Chain-of-Thought & 1.5-2x & +30-60\% \textit{nur bei Logik/Math} \\
Self-Consistency (5x) & 5x & +10-20\% \\
ReAct (Function Calling) & 2-5x & +50\% bei Multi-Step \\
\midrule
\end{tabularx}
\caption{Latenz vs. Qualität -- gemessen an gpt-4o-mini, Klassifikation + Logik Tasks}
\label{tab:latency-quality}
\end{table}

\realitycheck{Prompt-Länge = Kosten. 4.000 Token Prompt kostet 4x mehr als 1.000 Token -- bei oft marginalem Qualitätsgewinn. Miss \texttt{prompt\_efficiency\_score()}. Optimiere auf Tokens pro Qualitätspunkt.}

\subsection{Prompt Caching -- der einfachste Kostensenker}

System-Prompts + Few-Shot Examples ändern sich selten. Nutze Provider-Caching:

\begin{itemize}
    \item \textbf{OpenAI}: Automatisch ab 1.024 Tokens -- wiederholte Prefix-Tokens werden gecached (50\% Rabatt auf cached Tokens)
    \item \textbf{Anthropic}: Explizites \texttt{cache\_control: \{\textbackslash"type\textbackslash": \textbackslash"ephemeral\textbackslash"\}} auf System-Prompt + Examples
    \item \textbf{Template Pre-Rendering}: Statische Teile (System + Few-Shot) einmal rendern, nur User-Input dynamisch
\end{itemize}

\realitycheck{Ein guter Cache spart 40-60\% API-Kosten. Implementiere semantisches Caching (Embedding-Similarity auf User-Input) \textit{bevor} du über Model Routing nachdenkst. Größter Hebel, geringster Aufwand.}

\subsection{Token-Effizienz messen}

\begin{lstlisting}[language=Python, caption={Prompt Efficiency Score: Qualität pro Token}]
import tiktoken

def prompt_efficiency_score(prompt: str, output: str, task_quality: float) -> float:
    """Höher = besser. Viele Tokens für wenig Qualität = schlecht."""
    enc = tiktoken.get_encoding("o200k_base")  # gpt-4o tokenizer
    prompt_tokens = len(enc.encode(prompt))
    return task_quality / max(prompt_tokens, 1) * 1000

# Beispiel aus meiner Prod-Messung:
score_long  = prompt_efficiency_score(prompts["verbose"], outputs["verbose"], 0.95)  # 1.2
score_short = prompt_efficiency_score(prompts["compact"], outputs["compact"], 0.92)  # 3.8
# score_short > score_long: 3% weniger Qualität, 68% weniger Tokens
# Entscheidung: compact Version deployen
\end{lstlisting}

% ============================================================
% SICHERHEIT -- PROMPT INJECTION (VORWÄRTSVERWEIS)
% ============================================================
\section{Sicherheit -- Prompt Injection: die \#1 Produktionslücke}

\realitycheck{Prompt Injection ist \textbf{kein} theoretisches Problem. Jeder User-Input \textit{kann} bösartige Anweisungen enthalten. "Ignore previous instructions" ist der einfachste Angriff. Die goldene Regel: \textbf{User-Input niemals in den System-Prompt interpolieren.}}

Die vollständige Defense-in-Depth-Strategie (System-Prompt Rules, Input-Isolation via XML-Tags, Output-Validation via Pydantic/Zod, Moderation API, Rate-Limiting, Audit-Log) findest du in \textbf{Kapitel 16 (Security)}. Dort auch: Indirect Injection via RAG-Chunks, Multimodal Injection via Bilder, und EU-DSGVO-konforme Deployment-Patterns.

Hier nur das Minimum für Produktion:
\begin{enumerate}
    \item \textbf{System-Prompt}: Explizite Anti-Injection-Regeln + XML-Tag-Isolation
    \item \textbf{Input}: HTML-Escape + XML-Tag-Wrapping vor LLM-Call
    \item \textbf{Output}: Pydantic/Zod Schema-Validation (fängt Halluzination + Injection)
    \item \textbf{Moderation API}: Auf User-Input \textbf{und} LLM-Output
    \item \textbf{Audit-Log}: Jeder Prompt + Response für Forensik
\end{enumerate}

\realitycheck{Defense-in-Depth ist nicht optional: (1) System-Prompt Rules, (2) Input Sanitization/XML-Tags, (3) Output Schema Validation, (4) Moderation API, (5) Rate Limiting, (6) Audit Log jeder Prompt+Response. Das Minimum für Produktion: 1-3.}

% ============================================================
% ZUSAMMENFASSUNG -- KOMPAKT
% ============================================================
\section{Zusammenfassung}

\begin{itemize}
    \item Prompt Engineering = Hauptprogrammierschnittstelle des AI Engineers. Behandle es wie Code.
    \item Few-Shot (3-5 Edge-Cases) ist Produktionsstandard. Zero-Shot = Baseline.
    \item CoT nur für Logik/Math/Planning. ReAct = Agent-Basis (nutze Function Calling, nicht Text-Parsing).
    \item Self-Consistency (5x Cost) nur für Mission-Critical.
    \item Temperature 0.0-0.1 in Produktion. Seed für Reproduzierbarkeit (best effort).
    \item Prompt Injection: User-Input \textbf{niemals} in System-Prompt. XML-Isolation + Pydantic Validation + Moderation API.
    \item Prompts in Jinja2-Dateien, versioniert (Git-Tags), evaluiert (Golden Set), gecached (Provider + Semantic).
    \item Miss \texttt{prompt\_efficiency\_score()}. Kürzer = billiger + oft besser.
    \item Cache first, route later. Semantic Cache spart 40-60\% Kosten. Model Routing ist Optimierung \#2.
\end{itemize}

% ============================================================
% MERKE-BOX -- CHIP STYLE
% ============================================================



\par\mbox{}\par
\begin{merkeEnv}
\begin{itemize}
    \item \textbf{Ein Prompt = eine Aufgabe}. Multi-Task verwässert Qualität.
    \item \textbf{Temperature $>$ 0.3 in Produktion} = Fahrlässigkeit. Deterministisch = verlässlich.
    \item \textbf{System-Prompt ist heilig}. Selten ändern, immer mit Golden Set evaluieren.
    \item \textbf{Prompt Injection ist real}. Jeder User-Input ist potenzieller Angriff. Isolation ist Pflicht.
    \item \textbf{Kürzer ist besser}. Instruktionen kompakt, Examples präzise (Edge-Cases), Kontext minimal.
    \item \textbf{Golden Set vor zweitem Prompt}. Ohne automatisierte Eval fliegst du blind.
    \item \textbf{Cache first, route later}. Semantic Cache spart 40-60\% Kosten. Model Routing ist Optimierung \#2.
\end{itemize}
\end{merkeEnv}

% ============================================================
% PRAXISPROJEKT -- REALISTISCH
% ============================================================
\section{Praxisprojekt: Prompt-Engineering-Suite für Support-Classifier}

\autornotiz{
Das ist exakt das System, das ich 2023 für einen E-Commerce-Kunden gebaut habe.
20 Tickets/Min, 99.9\% Uptime, \$150/Monat API-Kosten.
Die Injection-Versuche kamen \textit{täglich} -- Competitor-Bots, neugierige User, Pen-Tests.
Die XML-Isolation + Pydantic-Validierung haben \textbf{alle} gefangen.
}

\subsection{Aufgabe: Produktionsreife Prompt-Engineering-Suite}

Baue ein vollständiges System für E-Commerce-Support-Ticket-Klassifikation.

\textbf{Anforderungen:}
\begin{itemize}[leftmargin=*]
    \item Prompt-Registry mit Jinja2 (mindestens 3 Versionen: v1 Zero-Shot, v2 Few-Shot, v3 Few-Shot + CoT für komplexe Cases)
    \item Alle 5 Techniken als wählbare Strategie (Enum: \texttt{ZERO\_SHOT, FEW\_SHOT, COT, REACT, SELF\_CONSISTENCY})
    \item Token-Counting \textbf{vor} API-Call (Budget-Enforcement: max 4k Input Tokens)
    \item Prompt-Injection-Schutz: XML-Isolation + Sanitization + Pydantic Validation + Moderation API
    \item Golden Set: 50 Test-Cases (30 normal, 10 Edge-Cases, 10 Injection-Versuche)
    \item CI-Pipeline: \texttt{pytest} läuft Golden Set bei jedem Prompt-Change. Fail bei Accuracy $<$ 90\% oder Injection-Pass-Rate $>$ 0\%
    \item Logging: Prompt-Hash, Version, Tokens, Latenz, Kosten, Validation-Status pro Request
\end{itemize}

\textbf{Testdaten:} 50 Tickets (siehe \texttt{tests/golden\_set.json} im Repo).

\textbf{Erfolgskriterium:} Classifier $>$ 90\% Accuracy. \textbf{100\%} Injection-Versuche erkannt \& abgewehrt. Token-Kosten pro 1k Requests dokumentiert. CI grün.

% ============================================================
% INTERVIEW FRAGEN -- WAS ICH FRAGE
% ============================================================
\section{Meine Interview-Fragen für AI Engineers}

\begin{enumerate}[leftmargin=*]
    \item \textbf{Frage:} "PM sagt: 'Der neue Prompt fühlt sich besser an.' Wie baust du Bauchgefühl in CI/CD um?"
    
    \textbf{Was ich suche:} Golden Set (50-100 Cases). Deterministisches (JSON) mit hartem Code testen. Semantische Qualität per LLM-as-a-Judge mit \textit{fixer} Rubrik. Keine Vibe-Checks in Produktion. Automatisierte Eval-Pipeline bei \textit{jedem} Prompt-Change.
    
    \item \textbf{Frage:} "Kernanwendung fällt wegen OpenAI-Downtime aus. Wie designst du drumherum?"
    
    \textbf{Was ich suche:} Abstraktions-Layer (LiteLLM / Custom Gateway) mit Fallback-Routen. Fällt Provider A mit 5xx aus $\rightarrow$ transparentes Routing auf Provider B. Health-Checks pro Model. Circuit Breaker.
    
    \item \textbf{Frage:} "Agent ruft Wetter-API falsch auf und hängt in Endlos-Loop. Prävention?"
    
    \textbf{Was ich suche:} (1) Harte Max-Iterationen im Agent-Loop (z.B. 10). (2) Tool-Parameter strikt per Pydantic validieren \textit{vor} Call. (3) Fehlercode als System-Message zurück an LLM ("Tool failed: invalid param X. Fix and retry."). (4) Timeout pro Tool-Call.
    
    \item \textbf{Frage:} "Wie verhinderst du, dass ein Prompt-Change die Kosten verdreifacht?"
    
    \textbf{Was ich suche:} Token-Budget-Check vor Call (max Input Tokens). Prompt-Hash in Traces. Alert bei Kosten-Anomalie (>2x Baseline). Prompt-Efficiency-Score in CI. Model Routing: Cheap Model Default, Expensive nur bei Low Confidence / Complex Query.
\end{enumerate}

% ============================================================
% WEITERFÜHRENDE RESSOURCEN -- KURATIERT
% ============================================================
\section{Weiterführende Ressourcen -- was ich wirklich lese}

\subsection{Dokumentationen (Primary Sources)}
\begin{itemize}[leftmargin=*]
    \item OpenAI Prompt Engineering Guide -- \href{https://platform.openai.com/docs/guides/prompt-engineering}{platform.openai.com/docs/guides/prompt-engineering}
    \item Anthropic Prompt Engineering -- \href{https://docs.anthropic.com/en/docs/build-with-claude/prompt-engineering}{docs.anthropic.com}
    \item Google Prompting Guide -- \href{https://ai.google.dev/gemini-api/docs/prompting}{ai.google.dev/gemini-api/docs/prompting}
\end{itemize}

\subsection{Tools, die ich nutze}
\begin{itemize}[leftmargin=*]
    \item \textbf{Langfuse} -- \href{https://langfuse.com}{langfuse.com} -- Open-Source Prompt Management, Tracing, Evals. Self-hostable.
    \item \textbf{Jinja2} -- \href{https://jinja.palletsprojects.com}{jinja.palletsprojects.com} -- Template Engine, Standard.
    \item \textbf{Weights \& Biases Prompts} -- \href{https://wandb.ai}{wandb.ai} -- Prompt Versionierung + Eval UI.
    \item \textbf{LangChain Prompt Hub} -- \href{https://smith.langchain.com/hub}{smith.langchain.com/hub} -- Geteilte Templates (inspirieren, nicht kopieren).
\end{itemize}

\subsection{Papers -- die Fundamente}
\begin{itemize}[leftmargin=*]
    \item Wei et al. (2022): "Chain-of-Thought Prompting Elicits Reasoning in Large Language Models" -- CoT-Grundlagen
    \item Yao et al. (2023): "ReAct: Synergizing Reasoning and Acting in Language Models" -- ReAct Paper
    \item Wang et al. (2023): "Self-Consistency Improves Chain of Thought Reasoning in Language Models"
    \item Schulhoff et al. (2024): "Ignore This Title and HackAPrompt: Exposing Systemic Vulnerabilities of LLMs" -- Injection Taxonomie
\end{itemize}

\autornotiz{
Nächstes Kapitel: Evaluation und Metriken -- der Qualitäts-Sicherungs-Riegel.
Da lernst du: Golden Sets (50-100 Cases). LLM-as-a-Judge mit Multi-Judge gegen Bias. Faithfulness, Answer Relevance, Context Precision. CI-Gate: kein Prompt-Change ohne Green-Build. Cost/Latency/Quality Triangle. Tracing (Langfuse) ab Tag 1. Ohne Eval fliegst du blind.
}

\textbf{Nächster Schritt:} Mit Prompt-Disziplin etabliert geht es weiter zu Evaluation -- wie du misst, ob dein System wirklich gut arbeitet, bevor es in Produktion geht.
\endinput

\clearpage
